\ifdefined\gkver\else\def\gkver{student}\fi
\documentclass[\gkver]{gaokaozhenti}
\begin{document}

\chapter{2015年10月浙江}

\begin{enumerate}
\item
%% number: 1
%% paperName: 2015年10月浙江选考第 1 题 3 分
%% typeId: 1
%% type: 单选题
%% chapter: 运动和力
%% point: 力学单位制
%% method: 无
%% score: 3
%% degree: 950
%% duplicateId: 0
%% body:
kg 和 s 是国际单位制两个基本单位的符号，这两个基本单位对应的物理量是\xzanswer{A}
\fourchoices[answer=A]
{质量和时间}
{质量和位移}
{重力和时间}
{重力和位移}
%% answer: A
%% memo:
\memoanswer{
kg 是质量的单位，s 是时间的单位；位移的单位是 m，重力的单位是 N，故 A 正确，B、C、D 错误。
}

\item
%% number: 2
%% paperName: 2015年10月浙江选考第 2 题 3 分
%% typeId: 1
%% type: 单选题
%% chapter: 直线运动
%% point: 匀速直线运动的 v-t 图像
%% method: 图像法
%% score: 3
%% degree: 900
%% duplicateId: 0
%% body:
下列 $v-t$ 图像中，表示物体做匀速直线运动的是\xzanswer{D}
\fourchoices[answer=D, ispicture=true, h=2.7cm]
{figs/02a.pdf}
{figs/02b.pdf}
{figs/02c.pdf}
{figs/02d.pdf}
%% answer: D
%% memo:
\memoanswer{
物体做匀速直线运动时速度保持不变，其 $v-t$ 图线与时间轴平行，故 D 正确。
}

\item
%% number: 3
%% paperName: 2015年10月浙江选考第 3 题 3 分
%% typeId: 1
%% type: 单选题
%% chapter: 直线运动
%% point: 参考系
%% method: 无
%% score: 3
%% degree: 900
%% duplicateId: 0
%% body:
2015 年 9 月 3 日，纪念中国人民抗日战争暨世界反法西斯战争胜利 70 周年阅兵式在天安门广场举行。如图所示，七架战机保持“固定队列”在天安门广场上空飞过。下列说法正确的是\xzanswer{A}
\onepicture[w=4.8cm]{figs/03.pdf}
\fourchoices[answer=A]
{以某飞机为参考系，其他飞机是静止的}
{以飞行员为参考系，广场上的观众是静止的}
{以某飞行员为参考系，其他飞行员是运动的}
{以广场上的观众为参考系，飞机是竖直向上运动的}
%% answer: A
%% memo:
\memoanswer{
七架战机保持“固定队列”飞行，以某飞机为参考系，其他飞机的位置不变，是静止的，故 A 正确；以飞行员为参考系，广场上的观众向后运动，故 B 错误；以某飞行员为参考系，其他飞行员是静止的，故 C 错误；以广场上的观众为参考系，飞机在水平方向上运动，故 D 错误。
}

\item
%% number: 4
%% paperName: 2015年10月浙江选考第 4 题 3 分
%% typeId: 1
%% type: 单选题
%% chapter: 直线运动
%% point: 匀速直线运动
%% method: 无
%% score: 3
%% degree: 850
%% duplicateId: 0
%% body:
小李乘坐高铁，当他所在的车厢刚要进隧道时，看到车厢内显示屏上的示数为 $216\Ukmh$，他立即观察手表秒针走动，经过 $20\Us$ 车厢出了隧道，则该隧道的长度约为\xzanswer{B}
\onepicture[w=4.8cm]{figs/04.pdf}
\fourchoices[answer=B]
{$600\Um$}
{$1200\Um$}
{$2160\Um$}
{$4320\Um$}
%% answer: B
%% memo:
\memoanswer{
由匀速直线运动公式得 $x=vt=\frac{216}{3.6}\times20\Um=1200\Um$，故 B 正确。
}

\item
%% number: 5
%% paperName: 2015年10月浙江选考第 5 题 3 分
%% typeId: 1
%% type: 单选题
%% chapter: 功和能
%% point: 能量守恒定律
%% method: 无
%% score: 3
%% degree: 900
%% duplicateId: 0
%% body:
画作《瀑布》如图所示。有人对此画作了如下解读：水流从高处倾泻而下，推动水轮机发电，又顺着水渠流动，回到瀑布上方，然后再次倾泻而下，如此自动地周而复始。这一解读违背了\xzanswer{D}
\onepicture[h=4.5cm]{figs/05.pdf}
\fourchoices[answer=D]
{库仑定律}
{欧姆定律}
{电荷守恒定律}
{能量守恒定律}
%% answer: D
%% memo:
\memoanswer{
水流从高处倾泻而下，推动水轮机发电，重力势能转化为电能；水流不可能自发地流回高处，这一解读违背了能量守恒定律，故 D 正确。
}

\item
%% number: 6
%% paperName: 2015年10月浙江选考第 6 题 3 分
%% typeId: 1
%% type: 单选题
%% chapter: 直线运动
%% point: 平均速度
%% method: 无
%% score: 3
%% degree: 850
%% duplicateId: 0
%% body:
如图所示，一女同学穿着轮滑鞋以一定的速度俯身“滑入”静止汽车的车底，她用 $15\Us$ 穿越了 20 辆汽车底部后“滑出”，位移为 $58\Um$。假设她的运动可视为匀变速直线运动，从上述数据可以确定\xzanswer{B}
\onepicture[w=5.0cm]{figs/06.pdf}
\fourchoices[answer=B]
{她在车底运动时的加速度}
{她在车底运动时的平均速度}
{她刚“滑入”车底时的速度}
{她刚“滑出”车底时的速度}
%% answer: B
%% memo:
\memoanswer{
由题意知，该同学运动的位移和时间已知，能通过计算确定其平均速度，故 B 正确；由于初速度未知，加速度及初、末速度均无法确定，故 A、C、D 错误。
}

\item
%% number: 7
%% paperName: 2015年10月浙江选考第 7 题 3 分
%% typeId: 1
%% type: 单选题
%% chapter: 曲线运动
%% point: 人造地球卫星
%% method: 无
%% score: 3
%% degree: 750
%% duplicateId: 0
%% body:
2015 年 9 月 20 日“长征六号”火箭搭载 20 颗小卫星成功发射。在多星分离时，小卫星分别在高度不同的三层轨道被依次释放。假设释放后的小卫星均做匀速圆周运动，则下列说法正确的是\xzanswer{C}
\onepicture[h=4.2cm]{figs/07.pdf}
\fourchoices[answer=C]
{20 颗小卫星的轨道半径均相同}
{20 颗小卫星的线速度大小均相同}
{同一圆轨道上的小卫星的周期均相同}
{不同圆轨道上的小卫星的角速度均相同}
%% answer: C
%% memo:
\memoanswer{
小卫星分别在高度不同的三层轨道被释放，不同位置的小卫星轨道半径不同，故 A 错误。由牛顿第二定律得 $G\frac{Mm}{r^{2}}=m\frac{v^{2}}{r}$，解得 $v=\sqrt{\frac{GM}{r}}$，20 颗小卫星的轨道半径 $r$ 不同，则线速度 $v$ 也都不同，故 B 错误。由 $G\frac{Mm}{r^{2}}=mr\left(\frac{2\pi}{T}\right)^{2}$ 得 $T=2\pi\sqrt{\frac{r^{3}}{GM}}$，同一轨道上的小卫星轨道半径 $r$ 相同，则周期 $T$ 均相同，故 C 正确。由 $G\frac{Mm}{r^{2}}=mr\omega^{2}$ 得 $\omega=\sqrt{\frac{GM}{r^{3}}}$，不同轨道上的小卫星轨道半径不同，则角速度不同，故 D 错误。
}

\item
%% number: 8
%% paperName: 2015年10月浙江选考第 8 题 3 分
%% typeId: 1
%% type: 单选题
%% chapter: 曲线运动
%% point: 竖直平面圆周运动
%% method: 机械能守恒
%% score: 3
%% degree: 700
%% duplicateId: 0
%% body:
质量为 $30\Ukg$ 的小孩坐在秋千板上，秋千板离系绳子的横梁的距离是 $2.5\Um$。小孩的父亲将秋千板从最低点拉起 $1.25\Um$ 高度后由静止释放，小孩沿圆弧运动至最低点时，她对秋千板的压力约为\xzanswer{C}
\fourchoices[answer=C]
{0}
{$200\UN$}
{$600\UN$}
{$1000\UN$}
%% answer: C
%% memo:
\memoanswer{
由机械能守恒得 $mgh=\frac{1}{2}mv^{2}$，解得 $v=\sqrt{2gh}=\sqrt{2\times10\times1.25}\Ums=5\Ums$。以小孩为研究对象，设秋千板对小孩的支持力为 $F_{N}$，由牛顿第二定律得 $F_{N}-mg=m\frac{v^{2}}{r}$，解得 $F_{N}=m\left(g+\frac{v^{2}}{r}\right)=30\times\left(10+\frac{5^{2}}{2.5}\right)\UN=600\UN$。由牛顿第三定律知，小孩对秋千板的压力为 $600\UN$，故 C 正确。
}

\item
%% number: 9
%% paperName: 2015年10月浙江选考第 9 题 3 分
%% typeId: 1
%% type: 单选题
%% chapter: 磁场
%% point: 霍尔效应
%% method: 控制变量法
%% score: 3
%% degree: 700
%% duplicateId: 0
%% body:
小张在探究磁场对电流作用的实验中，将直导线换作导体板，如图所示，发现在 a、b 两点之间存在电压 $U_{ab}$。进一步实验结果如下表：
\onepicture[w=5.5cm]{figs/09.pdf}

\begin{center}
\begin{tabular}{|c|c|c|}
\hline
电流 & 磁感应强度 & 电压 $U_{ab}$ \\ \hline
I & B & U \\ \hline
I & 2B & 2U \\ \hline
I & 3B & 3U \\ \hline
2I & B & 2U \\ \hline
3I & B & 3U \\ \hline
\end{tabular}
\end{center}

由表中结果可知电压 $U_{ab}$\xzanswer{C}
\fourchoices[answer=C]
{与电流无关}
{与磁感应强度无关}
{与电流可能成正比}
{与磁感应强度可能成反比}
%% answer: C
%% memo:
\memoanswer{
由表格数据可知，当控制磁感应强度 $B$ 相同时，电压 $U$ 与电流 $I$ 成正比，故 A 错误、C 正确；当控制电流 $I$ 相同时，电压 $U$ 与磁感应强度 $B$ 成正比，故 B、D 错误。
}

\item
%% number: 10
%% paperName: 2015年10月浙江选考第 10 题 3 分
%% typeId: 1
%% type: 单选题
%% chapter: 电路
%% point: 电功电功率
%% method: 无
%% score: 3
%% degree: 750
%% duplicateId: 0
%% body:
如图所示是我国某 10 兆瓦（1 兆瓦 $=10^{6}\UW$）光伏发电站，投入使用后每年可减少排放近万吨二氧化碳。已知该地区每年能正常发电的时间约为 $1200\Uh$，则该电站年发电量约为\xzanswer{B}
\onepicture[w=6.5cm]{figs/10.pdf}
\fourchoices[answer=B]
{$1.2\times10^{6}\UkWh$}
{$1.2\times10^{7}\UkWh$}
{$1.2\times10^{9}\UkWh$}
{$1.2\times10^{10}\UkWh$}
%% answer: B
%% memo:
\memoanswer{
由 $W=Pt$ 知，发电站的年发电量 $E=Pt=10\times10^{6}\UW\times1200\Uh=10\times10^{3}\UkW\times1200\Uh=1.2\times10^{7}\UkWh$，故 B 正确。
}

\item
%% number: 11
%% paperName: 2015年10月浙江选考第 11 题 3 分
%% typeId: 1
%% type: 单选题
%% chapter: 电场
%% point: 库仑定律
%% method: 受力分析
%% score: 3
%% degree: 650
%% duplicateId: 0
%% body:
如图所示，一质量为 $m$、电荷量为 $Q$ 的小球 A 系在长为 $l$ 的绝缘轻绳下端，另一电荷量也为 $Q$ 的小球 B 位于悬挂点的正下方（A、B 均视为点电荷），轻绳与竖直方向成 $30^{\circ}$ 角，小球 A、B 静止于同一高度。已知重力加速度为 $g$，静电力常量为 $k$，则两球间的静电力为\xzanswer{A}
\onepicture[h=4.2cm]{figs/11.pdf}
\fourchoices[answer=A]
{$\frac{4kQ^{2}}{l^{2}}$}
{$\frac{kQ^{2}}{l^{2}}$}
{$mg$}
{$\sqrt{3}mg$}
%% answer: A
%% memo:
\memoanswer{
带电小球 A 处于平衡状态，受到库仑力 $F$、重力 $mg$ 以及绳子的拉力 $T$ 的作用，其合力为零。由受力分析得 $F=mg\tan30^{\circ}=\frac{\sqrt{3}}{3}mg$；由库仑定律有 $F=k\frac{Q^{2}}{r^{2}}$，因 $r=l\sin30^{\circ}=\frac{l}{2}$，所以 $F=\frac{4kQ^{2}}{l^{2}}$，故 A 正确，B、C、D 错误。
}

\item
%% number: 12
%% paperName: 2015年10月浙江选考第 12 题 3 分
%% typeId: 1
%% type: 单选题
%% chapter: 功和能
%% point: 功和功率
%% method: 无
%% score: 3
%% degree: 700
%% duplicateId: 0
%% body:
快艇在运动中受到的阻力与速度平方成正比（即 $F_{f}=kv^{2}$）。若油箱中有 $20\UL$ 燃油，当快艇以 $10\Ums$ 匀速行驶时，还能行驶 $40\Ukm$，假设快艇发动机的效率保持不变，则快艇以 $20\Ums$ 匀速行驶时，还能行驶\xzanswer{C}
\fourchoices[answer=C]
{$80\Ukm$}
{$40\Ukm$}
{$10\Ukm$}
{$5\Ukm$}
%% answer: C
%% memo:
\memoanswer{
快艇匀速行驶时，牵引力与阻力平衡，即 $F=F_{f}=kv^{2}$。油箱中燃油相同且发动机效率保持不变，说明两种情况下提供的能量相等，克服阻力做的功相等。由 $W=Fs$ 得 $kv_{1}^{2}s_{1}=kv_{2}^{2}s_{2}$，又 $v_{2}=2v_{1}$，所以 $s_{2}=\frac{v_{1}^{2}}{v_{2}^{2}}s_{1}=\frac{1}{4}\times40\Ukm=10\Ukm$，故 C 正确。
}

\item
%% number: 13
%% paperName: 2015年10月浙江选考第 13 题 3 分
%% typeId: 1
%% type: 单选题
%% chapter: 力物体的平衡
%% point: 摩擦力
%% method: 图像法
%% score: 3
%% degree: 750
%% duplicateId: 0
%% body:
将质量为 $1.0\Ukg$ 的木块放在水平长木板上，用力沿水平方向拉木块，拉力从 0 开始逐渐增大，木块先静止后相对木板运动。用力传感器采集木块受到的拉力和摩擦力的大小，并用计算机绘制出摩擦力大小 $F_{f}$ 随拉力大小 $F$ 变化的图像，如图所示。木块与长木板间的动摩擦因数为\xzanswer{A}
\onepicture[w=4.5cm]{figs/13.pdf}
\fourchoices[answer=A]
{0.3}
{0.5}
{0.6}
{1.0}
%% answer: A
%% memo:
\memoanswer{
由题图知，木块运动后所受的滑动摩擦力为 $F_{f}=3\UN$。由 $F_{f}=\mu F_{N}=\mu mg$，解得 $\mu=\frac{F_{f}}{mg}=\frac{3}{1.0\times10}=0.3$，故 A 正确。
}

\item
%% number: 14
%% paperName: 2015年10月浙江选考第 14 题 2 分
%% typeId: 2
%% type: 多选题
%% chapter: 原子和原子核
%% point: 氢原子能级跃迁
%% method: 无
%% score: 2
%% degree: 650
%% duplicateId: 0
%% body:
【加试题】氢原子从 $n=6$ 跃迁到 $n=2$ 能级时辐射出频率为 $\nu_{1}$ 的光子，从 $n=5$ 跃迁到 $n=2$ 能级时辐射出频率为 $\nu_{2}$ 的光子。下列说法正确的是\xzanswer{ABD}
\fourchoices[answer=ABD]
{频率为 $\nu_{1}$ 的光子的能量较大}
{频率为 $\nu_{1}$ 的光子的动量较大}
{做双缝干涉实验时，频率为 $\nu_{1}$ 的光产生的条纹间距较大}
{做光电效应实验时，频率为 $\nu_{1}$ 的光产生的光电子的最大初动能较大}
%% answer: ABD
%% memo:
\memoanswer{
氢原子从高能级跃迁到低能级会辐射出光子，有 $h\nu=E_{\text{初}}-E_{\text{末}}$，可知频率为 $\nu_{1}$ 的光子能量较大，$\nu_{1}>\nu_{2}$，故 A 正确。由 $\lambda=\frac{h}{p}$、$\lambda=\frac{c}{\nu}$ 得光子的动量 $p=\frac{h\nu}{c}$，频率为 $\nu_{1}$ 的光子动量较大，故 B 正确。在双缝干涉中，相邻明条纹（或暗条纹）间距 $\Delta x=\frac{l}{d}\lambda=\frac{lc}{d\nu}$，频率为 $\nu_{1}$ 的光产生的条纹间距较小，故 C 错误。做光电效应实验时，$E_{km}=h\nu-W_{\text{逸}}$，频率为 $\nu_{1}$ 的光产生的光电子的最大初动能较大，故 D 正确。
}

\item
%% number: 15
%% paperName: 2015年10月浙江选考第 15 题 2 分
%% typeId: 2
%% type: 多选题
%% chapter: 机械波
%% point: 波的传播
%% method: 图像法
%% score: 2
%% degree: 550
%% duplicateId: 0
%% body:
【加试题】沿 x 轴正方向传播的简谐横波，振幅为 $5\Ucm$，质点 A、B 的平衡位置离 O 点的距离分别为 $1.8\Um$ 和 $3.0\Um$。$t=0$ 时的波形如图所示，$t=0.6\Us$ 时质点 A 恰好第二次到达波峰。下列说法正确的是\xzanswer{BD}
\onepicture[w=6.5cm]{figs/15.pdf}
\fourchoices[answer=BD]
{波速为 $3\Ums$}
{频率为 $2\UHz$}
{质点 B 在 $0\sim0.6\Us$ 内路程为 $25\Ucm$}
{质点 B 在 $t=0.6\Us$ 时沿 y 轴正方向运动}
%% answer: BD
%% memo:
\memoanswer{
由图像可知，波长 $\lambda=4\Um$。波沿 x 轴正方向传播，A 点此时向上振动，当 $x=1\Um$ 处的振动传到 $x=1.8\Um$ 处的 A 点时，A 点第一次到达波峰；$t=0.6\Us$ 时质点 A 恰好第二次到达波峰，则 $0.6=\frac{1.8-1}{v}+T$，又 $v=\frac{\lambda}{T}=\frac{4}{T}$，联立解得 $v=8\Ums$，$T=0.5\Us$，故 A 错误。频率 $f=\frac{1}{T}=2\UHz$，故 B 正确。$0.6\Us=1\frac{1}{5}T<1\frac{1}{4}T$，而 B 点在 $1\frac{1}{4}T$ 内的路程为 $s=1\frac{1}{4}\times4A=\frac{5}{4}\times4\times5\Ucm=25\Ucm$，所以质点 B 在 $0\sim0.6\Us$ 内路程小于 $25\Ucm$，故 C 错误。由图像可知，B 在 $t=0$ 时刻在波谷处向上振动，则 $t=0.6\Us=1\frac{1}{5}T<1\frac{1}{4}T$ 时在平衡位置下方沿 y 轴正方向运动，故 D 正确。
}

\item
%% number: 16
%% paperName: 2015年10月浙江选考第 16 题 2 分
%% typeId: 2
%% type: 多选题
%% chapter: 光学
%% point: 光的折射
%% method: 光路分析
%% score: 2
%% degree: 450
%% duplicateId: 0
%% body:
【加试题】为了从坦克内部观察外部目标，在厚度为 $20\Ucm$ 的坦克壁上开了一个直径为 $12\Ucm$ 的孔，若在孔内分别安装由同一材料制成的如图所示的三块玻璃，其中两块玻璃的厚度相同。坦克内的人在同一位置通过玻璃能看到的外界的角度范围是\xzanswer{AD}
\threepicture[showlabel=false, h=3.0cm]
{figs/16a.pdf}
{figs/16b.pdf}
{figs/16c.pdf}
\fourchoices[answer=AD]
{图1的大于图2的}
{图1的小于图3的}
{图2的小于图3的}
{图2的等于图3的}
%% answer: AD
%% memo:
\memoanswer{
光线穿过玻璃砖后出射光线和入射光线平行，光路图如图所示。

	\threepicture[showlabel=false, h=3.0cm]{figs/16_an1.pdf}{figs/16_an2.pdf}{figs/16_an3.pdf}

	图1中光线在玻璃砖中的入射角大于图2、图3中玻璃砖中的入射角，由折射定律，射入空气的折射角图1最大；由光的可逆原理，图1观察到的角度范围最大。由图可知，出射光线和入射光线平行，在图2和图3中，由几何关系知，两光线在玻璃砖中的入射角相等，则在空气中的折射角也相等，即图2和图3观察的范围相同，故 A、D 正确，B、C 错误。
}

\item
%% number: 17
%% paperName: 2015年10月浙江选考第 17 题 5 分
%% typeId: 4
%% type: 实验题
%% chapter: 力物体的平衡
%% point: 探究求合力的方法
%% method: 等效替代法
%% score: 5
%% degree: 700
%% duplicateId: 0
%% body:
\begin{enumerate}
	\item
	在“探究求合力的方法”实验中，下列器材中必须要用的是（多选）。

	\fourchoices[answer=BD, ispicture=true, h=2.8cm]
	{figs/17a.pdf}
	{figs/17b.pdf}
	{figs/17c.pdf}
	{figs/17d.pdf}

	\item
	下列两个学生实验中，必须要测量质量、时间和位移的是\tkanswer{B}。选项 A：用打点计时器测速度；选项 B：探究加速度与力、质量的关系。
\end{enumerate}
%% answer:
\jdanswer{
\begin{enumerate}
	\item
	BD

	\item
	B

\end{enumerate}
}
%% memo:
\memoanswer{
（1）做“探究共点力合成的规律”实验，需要让两个力拉橡皮条和一个力拉橡皮条产生的作用效果相同，测出两个力的大小和方向以及一个力的大小和方向，用力的图示画出这三个力，用平行四边形作出两个力的合力的理论值，再和那一个力进行比较。所以需要的实验器材有：方木板（固定白纸）、白纸（记录方向画图）、刻度尺（选标度）、绳套（弹簧秤拉橡皮条）、弹簧测力计（测力的大小）、图钉（固定白纸）、三角板（画平行四边形）、橡皮条（让力产生相同的作用效果）。不需要秒表和打点计时器，故 B、D 正确，A、C 错误。

（2）用打点计时器测速度不需要测量质量；探究加速度与力、质量的关系的实验需要测量质量、时间和位移，故 B 正确。
}

\item
%% number: 18
%% paperName: 2015年10月浙江选考第 18 题 5 分
%% typeId: 4
%% type: 实验题
%% chapter: 电路
%% point: 测定电池的电动势和内阻
%% method: 图像法
%% score: 5
%% degree: 650
%% duplicateId: 0
%% body:
某实验小组做“测定电池的电动势和内阻”的实验，实验电路如图1所示。
\begin{enumerate}
	\item
	电流表示数如图2所示，则电流为\tkanswer{0.25}$\UA$。

	\twopicture[showlabel=false, align=right, wA=3.6cm, wB=5.6cm]
	{figs/18a.pdf}
	{figs/18b.pdf}

	\item
	根据实验数据获得 $U-I$ 图像如图3所示，则电动势 $E=$\tkanswer{1.43}$\UV$，内阻 $r=$\tkanswer{0.75}$\UO$。

	\onepicture[w=6.5cm]{figs/18c.pdf}
\end{enumerate}
%% answer:
\jdanswer{
\begin{enumerate}
	\item
	$0.25\pm0.01$

	\item
	$1.43\pm0.01$；$0.75\pm0.03$

\end{enumerate}
}
%% memo:
\memoanswer{
（1）电流表选择的是 $0\sim0.6\UA$ 量程，刻度盘中每一小格是 $0.02\UA$，故读数时不需要再估读，直接读数为 $0.25\UA$。

（2）由实验电路，根据闭合电路欧姆定律得 $U=E-Ir$，$U-I$ 图像斜率的绝对值代表内阻，由图3可知 $r=\left|\frac{1.43\UV-1.0\UV}{0.57\UA}\right|=0.75\UO$；$I=0$ 时，图线与 $U$ 轴的交点等于 $E$，所以 $E=1.43\UV$。
}

\item
%% number: 19
%% paperName: 2015年10月浙江选考第 19 题 9 分
%% typeId: 6
%% type: 计算题
%% chapter: 直线运动
%% point: 匀变速直线运动与平抛运动
%% method: 无
%% score: 9
%% degree: 700
%% duplicateId: 0
%% body:
如图1所示，饲养员对着长 $l=1.0\Um$ 的水平细长管的一端吹气，将位于吹气端口的质量 $m=0.02\Ukg$ 的注射器射到动物身上。注射器飞离长管末端的速度大小 $v=20\Ums$。可视为质点的注射器在长管内做匀变速直线运动，离开长管后做平抛运动，如图2所示。重力加速度 $g$ 取 $10\Umsq$。
\begin{enumerate}
	\item
	求注射器在长管内运动时的加速度大小；
	\item
	求注射器在长管内运动时受到的合力大小；
	\item
	若动物与长管末端的水平距离 $x=4.0\Um$，求注射器下降的高度 $h$。
\end{enumerate}
\twopicture[showlabel=false, wA=5.0cm, wB=7.0cm, align=right]
{figs/19a.pdf}
{figs/19b.pdf}
%% answer:
\jdanswer{
\begin{enumerate}
	\item
	$2.0\times10^{2}\Umsq$

	\item
	$4\UN$

	\item
	$0.2\Um$

\end{enumerate}
}
%% memo:
\memoanswer{
（1）由匀变速直线运动规律得 $v^{2}-0=2al$，解得 $a=\frac{v^{2}}{2l}=\frac{20^{2}}{2\times1.0}\Umsq=2.0\times10^{2}\Umsq$。

（2）由牛顿第二定律得 $F=ma$，解得 $F=0.02\times2.0\times10^{2}\UN=4\UN$。

（3）由平抛运动规律得 $x=vt$，解得 $t=\frac{x}{v}=\frac{4.0}{20}\Us=0.2\Us$；$h=\frac{1}{2}gt^{2}$，解得 $h=\frac{1}{2}\times10\times0.2^{2}\Um=0.2\Um$。
}

\item
%% number: 20
%% paperName: 2015年10月浙江选考第 20 题 12 分
%% typeId: 6
%% type: 计算题
%% chapter: 功和能
%% point: 动能定理
%% method: 无
%% score: 12
%% degree: 600
%% duplicateId: 0
%% body:
如图所示是公路上的“避险车道”，车道表面是粗糙的碎石，其作用是供下坡的汽车在刹车失灵的情况下避险。质量 $m=2.0\times10^{3}\Ukg$ 的汽车沿下坡行驶，当驾驶员发现刹车失灵的同时发动机失去动力，此时速度表示数 $v_{1}=36\Ukmh$，汽车继续沿下坡匀加速直行 $l=350\Um$、下降高度 $h=50\Um$ 时到达“避险车道”，此时速度表示数 $v_{2}=72\Ukmh$。重力加速度 $g$ 取 $10\Umsq$。
\begin{enumerate}
	\item
	求从发现刹车失灵至到达“避险车道”这一过程汽车动能的变化量；
	\item
	求汽车在下坡过程中所受的阻力；
	\item
	若“避险车道”与水平面间的夹角为 $17^{\circ}$，汽车在“避险车道”受到的阻力是在下坡公路上的 3 倍，求汽车在“避险车道”上运动的最大位移（$\sin17^{\circ}\approx0.3$）。
\end{enumerate}
\onepicture[align=right, w=7.0cm]{figs/20.pdf}
%% answer:
\jdanswer{
\begin{enumerate}
	\item
	$3.0\times10^{5}\UJ$

	\item
	$2\times10^{3}\UN$

	\item
	$33.3\Um$

\end{enumerate}
}
%% memo:
\memoanswer{
（1）由动能定理得 $\Delta E_{k}=\frac{1}{2}mv_{2}^{2}-\frac{1}{2}mv_{1}^{2}$，其中 $v_{1}=36\Ukmh=10\Ums$，$v_{2}=72\Ukmh=20\Ums$，解得 $\Delta E_{k}=3.0\times10^{5}\UJ$。

（2）由动能定理得 $mgh-F_{f}l=\frac{1}{2}mv_{2}^{2}-\frac{1}{2}mv_{1}^{2}$，解得 $F_{f}=\frac{\frac{1}{2}mv_{2}^{2}-\frac{1}{2}mv_{1}^{2}+mgh}{l}=2\times10^{3}\UN$。

（3）设汽车在“避险车道”上向上运动的最大位移为 $l'$，由动能定理得 $-(mg\sin17^{\circ}+3F_{f})l'=0-\frac{1}{2}mv_{2}^{2}$，解得 $l'=\frac{\frac{1}{2}mv_{2}^{2}}{mg\sin17^{\circ}+3F_{f}}=33.3\Um$。
}

\item
%% number: 21
%% paperName: 2015年10月浙江选考第 21 题 2 分
%% typeId: 4
%% type: 实验题
%% chapter: 电学实验
%% point: 变压器与电阻率测量
%% method: 无
%% score: 2
%% degree: 650
%% duplicateId: 0
%% body:
【加试题】
\begin{enumerate}
	\item
	在“探究变压器线圈两端的电压与匝数的关系”实验中，电源输出端与变压器线圈连接方式正确的是\tkanswer{图2}（填“图1”或“图2”）。变压器原、副线圈的匝数分别为 120 匝和 60 匝，测得的电压分别为 $8.2\UV$ 和 $3.6\UV$，据此可知电压比与匝数比不相等，主要原因是\tkanswer{漏磁、铁芯发热、导线发热}。

	\twopicture[showlabel=false, wA=2.8cm, wB=2.8cm, align=right]
	{figs/21a.pdf}
	{figs/21b.pdf}

	\item
	在“探究导体电阻与其影响因素的关系”实验中，某同学测量金属丝直径时，螺旋测微器示数如图3所示，读数是\tkanswer{0.638}$\Umm$；若要测量该金属丝的电阻率，还需要用到的测量仪器有\tkanswer{电压表、电流表、刻度尺}。

	\onepicture[align=right, w=5.5cm]{figs/21.pdf}
\end{enumerate}
%% answer:
\jdanswer{
\begin{enumerate}
	\item
	图2；漏磁、铁芯发热、导线发热

	\item
	$0.638\pm0.002$；电压表、电流表、刻度尺

\end{enumerate}
}
%% memo:
\memoanswer{
（1）变压器是靠电磁感应原理工作的，必须接交流电，故正确的是图2（原答案页作“图乙”，即图2）；理想变压器电压与匝数成正比，实际变压器不成正比的原因是由于漏磁、铁芯发热、导线发热等电能损失造成的。

（2）由螺旋测微器工作原理，图3读数为 $0.5\Umm+13.8\times0.01\Umm=0.638\Umm$。测量该金属丝的电阻率，由 $R=\rho\frac{l}{S}$ 得 $\rho=\frac{RS}{l}$，还需要测量金属丝的电阻 $R$ 和长度 $l$，故还需要用到的测量仪器有电压表、电流表、刻度尺。
}

\item
%% number: 22
%% paperName: 2015年10月浙江选考第 22 题 8 分
%% typeId: 6
%% type: 计算题
%% chapter: 电磁感应
%% point: 法拉第电磁感应定律
%% method: 无
%% score: 8
%% degree: 500
%% duplicateId: 0
%% body:
【加试题】如图1所示，质量 $m=3.0\times10^{-3}\Ukg$ 的“方框”形金属细框竖直放置在两水银槽中，方框的水平细杆 CD 长 $l=0.20\Um$，处于磁感应强度大小 $B_{1}=1.0\UT$、方向水平向右的匀强磁场中。有一匝数 $n=300$ 匝、面积 $S=0.01\Umq$ 的线圈通过开关 K 与两水银槽相连。线圈处于与线圈平面垂直的、沿竖直方向的匀强磁场中，其磁感应强度 $B_{2}$ 的大小随时间 $t$ 变化的关系如图2所示。重力加速度 $g$ 取 $10\Umsq$。
\begin{enumerate}
	\item
	求 $0\sim0.10\Us$ 线圈中的感应电动势大小；
	\item
	$t=0.22\Us$ 时闭合开关 K，若细杆 CD 所受安培力方向竖直向上，判断 CD 中的电流方向及磁感应强度 $B_{2}$ 的方向；
	\item
	$t=0.22\Us$ 时闭合开关 K，若安培力远大于重力，细框跳起的最大高度 $h=0.20\Um$，求通过细杆 CD 的电荷量。
\end{enumerate}
\twopicture[showlabel=false, wA=4.5cm, wB=4.8cm, align=right]
{figs/22a.pdf}
{figs/22b.pdf}
%% answer:
\jdanswer{
\begin{enumerate}
	\item
	$30\UV$

	\item
	CD 中的电流方向为 $C\to D$，$B_{2}$ 的方向向上

	\item
	$0.03\UC$

\end{enumerate}
}
%% memo:
\memoanswer{
（1）由法拉第电磁感应定律得 $E=n\frac{\Delta\Phi}{\Delta t}$，又 $\Delta\Phi=\Delta B_{2}S$，由图2得 $\frac{\Delta B_{2}}{\Delta t}=\frac{1.0}{0.10}\UT/\Us=10\UT/\Us$，解得 $E=30\UV$。

（2）利用左手定则分析知，CD 细杆中的电流方向为 $C\to D$；由感应电流的方向，利用右手螺旋定则分析知，磁感应强度 $B_{2}$ 的方向向上。

（3）由运动学公式得 $v^{2}=2gh$，解得 $v=\sqrt{2gh}$；由牛顿第二定律得 $F=ma$，又 $a=\frac{v-0}{\Delta t}$；细杆 CD 受到的安培力大小为 $F=B_{1}Il$；由电荷的定义式知 $\Delta Q=I\Delta t$，联立得 $\Delta Q=\frac{m\sqrt{2gh}}{B_{1}l}=0.03\UC$。
}

\item
%% number: 23
%% paperName: 2015年10月浙江选考第 23 题 12 分
%% typeId: 6
%% type: 计算题
%% chapter: 磁场
%% point: 带电粒子在磁场中的运动
%% method: 无
%% score: 12
%% degree: 400
%% duplicateId: 0
%% body:
【加试题】如图是水平放置的小型粒子加速器的原理示意图，区域Ⅰ和Ⅱ存在方向垂直纸面向里的匀强磁场 $B_{1}$ 和 $B_{2}$，长 $L=1.0\Um$ 的区域Ⅲ存在场强大小 $E=5.0\times10^{4}\UVm$、方向水平向右的匀强电场。区域Ⅲ中间上方有一离子源 S，水平向左发射动能 $E_{k0}=4.0\times10^{4}\UeV$ 的氘核，氘核最终从区域Ⅱ下方的 P 点水平射出。S、P 两点间的高度差 $h=0.10\Um$。

（氘核质量 $m=2\times1.67\times10^{-27}\Ukg$、电荷量 $q=1.60\times10^{-19}\UC$，$1\UeV=1.60\times10^{-19}\UJ$，$\sqrt{\frac{1.67\times10^{-27}}{1.60\times10^{-19}}}\approx1\times10^{-4}$）
\begin{enumerate}
	\item
	求氘核经过两次加速后从 P 点射出时的动能 $E_{k2}$；
	\item
	若 $B_{1}=1.0\UT$，要使氘核经过两次加速后从 P 点射出，求区域Ⅰ的最小宽度 $d$；
	\item
	若 $B_{1}=1.0\UT$，要使氘核经过两次加速后从 P 点射出，求区域Ⅱ的磁感应强度 $B_{2}$。
\end{enumerate}
\onepicture[align=right, w=9.0cm]{figs/23.pdf}
%% answer:
\jdanswer{
\begin{enumerate}
	\item
	$2.24\times10^{-14}\UJ$

	\item
	$0.06\Um$

	\item
	$1.2\UT$

\end{enumerate}
}
%% memo:
\memoanswer{
（1）由动能定理得 $W=E_{k2}-E_{k0}$，电场力做的功为 $W=qE\cdot2L$，代入数据解得 $E_{k2}=2.24\times10^{-14}\UJ$。（原答案页记为 $2.24\times10^{-34}\UJ$，显系排版或识别笔误，按计算应为 $2.24\times10^{-14}\UJ$。）

（2）由洛伦兹力提供向心力得 $qvB=m\frac{v^{2}}{r}$。第一次进入 $B_{1}$ 区域时，轨迹的半径为 $R_{0}=\frac{mv_{0}}{qB_{1}}=0.04\Um$；第二次进入 $B_{1}$ 区域，有 $\frac{1}{2}mv_{1}^{2}=E_{k0}+qEL$，$R_{2}=\frac{mv_{1}}{qB_{1}}=0.06\Um$，又 $d=R_{2}$，解得 $d=0.06\Um$，即区域Ⅰ的最小宽度为 $0.06\Um$。

（3）氘核运动轨迹如图所示。

	\onepicture[w=8.5cm]{figs/23_an.pdf}

	由图中几何关系可知 $2R_{2}=h+(2R_{1}-2R_{0})$，解得 $R_{1}=0.05\Um$；由 $R_{1}=\frac{mv_{1}}{qB_{2}}$ 得 $B_{2}=\frac{mv_{1}}{qR_{1}}=1.2\UT$。
}

\end{enumerate}

\end{document}
